%% ═══════════════════════════════════════════════════════════════════════════
\section{Register of Modelling Assumptions}
\label{app:assumptions}
%% ═══════════════════════════════════════════════════════════════════════════

A model is the set of commitments it makes, and those commitments are scattered through
Section~\ref{sec:Methods} at the point where each first becomes necessary. This appendix
collects them in one place so that a reader can weigh the whole set at once rather than
reconstruct it, and so that the ones this work could not test are as visible as the ones it
could.

Three statuses are used, and the distinction between the last two is the one that matters.
\emph{Tested} means the campaign produced evidence bearing on the assumption and
Section~\ref{sec:Results} reports it. \emph{Testable, untested} means an experiment within
reach of this design would bear on it and was not run, usually for want of events or time;
these are recoverable and Section~\ref{sec:future_work} carries them. \emph{Outside this
design} means no experiment this instrument can run would settle the question, so the
assumption is load-bearing and permanent until the design changes; these are the ones that
bound what the thesis may claim.

\subsection{The register}

Read down the status column first. It is the shorter question, and it decides how much
weight the rest of the row can carry.

\begin{longtable}{@{}p{0.5cm} p{4.1cm} p{2.0cm} p{7.6cm}@{}}
\caption[Register of modelling assumptions]{Every modelling assumption, where it is stated,
and what is known about it}
\label{tab:assumptions} \\
\toprule
& \textbf{Assumption} & \textbf{Status} & \textbf{If it fails} \\
\midrule
\endfirsthead
\multicolumn{4}{@{}l}{\small\itshape Table~\ref{tab:assumptions}, continued} \\
\toprule
& \textbf{Assumption} & \textbf{Status} & \textbf{If it fails} \\
\midrule
\endhead
\bottomrule
\endlastfoot

A1 & The six blocks are independent competing causes of rupture
(Assumption~\ref{as:noisyor}, Section~\ref{sec:model_ur})
& Tested
& Correlated blocks double-count a shared cause and the aggregate overstates urgency. The
direction is conservative: it errs towards alarm, never towards complacency.
Section~\ref{sec:res_armA} reports the observed inter-block correlations. \\
\addlinespace

A2 & The blocks are exchangeable up to their weights $\omega_m$
(Section~\ref{sec:model_ur})
& Tested, and false in detail
& $u_{\text{time}}$ is a frequentist probability and $u_{\text{cost}}$ an ordinal statement in
cardinal dress; no transfer function makes them equally informative. The aggregate is
therefore usable for ranking and for comparison against $U_d$, and is not a calibrated
probability. This is the reason no Brier score is reported on $U_r$ itself. \\
\addlinespace

A3 & The declaration is sincere (Assumption~\ref{as:sincerity},
Section~\ref{sec:model_adequacy})
& Outside this design
& Hidden risk stops measuring misperception and starts measuring what a declarant is willing
to write down. The two demand opposite interventions and are indistinguishable in this data.
Every behavioural reading in Section~\ref{sec:Results} rests on it. \\
\addlinespace

A4 & Language-model agents stand in for human declarants
(Section~\ref{sec:methods_arms})
& Outside this design
& Every statement about declaration behaviour applies to that population and not to
consultants. The mechanical findings are unaffected, since they concern code paths rather
than declarants. \\
\addlinespace

A5 & Lead times are lognormal by default, parameterised by moments
(Section~\ref{sec:model_mc})
& Testable, untested
& The tail of the completion-time distribution is misstated and $u_{\text{time}}$ with it.
The law is configurable per node, and normal-clipped and triangular alternatives are
implemented, so the sensitivity is measurable on frozen scenarios at low cost. \\
\addlinespace

A6 & Lead times are drawn independently across nodes
(Section~\ref{sec:model_mc})
& Outside this design
& A common shock, an energy crisis or a regional closure, correlates delays that the
simulation treats as independent, and the simulated completion distribution is too narrow.
Nothing in the campaign scenario could separate this from ordinary variance. \\
\addlinespace

A7 & The arc coefficients $\beta_{ji}$ and $\gamma_{ik}$ are known constants
(Section~\ref{sec:model_prop})
& Testable, untested
& Propagated urgency is misscaled. The coefficients are set by the scenario rather than
estimated from data, and estimating them from observed co-movement of node states is the
natural next study. \\
\addlinespace

A8 & Backup arcs carry neither propagation
(Section~\ref{sec:model_prop})
& Modelling choice, exercised
& A network with a real alternative supplier reads as more exposed than it is. The choice is
deliberate, since an unqualified backup is documentation rather than capacity, and the
campaign scenario exploits it on purpose. \\
\addlinespace

A9 & The Prospect Theory parameters transfer to supply chain coordinators,
$\lambda_{\text{under}} = 2.25$ and $\alpha = 0.88$
(Section~\ref{sec:model_adequacy})
& Testable, untested
& The adequacy score is miscalibrated as a governance instrument, though its ordering is
unaffected, since $A$ is monotone in the gap whatever the parameters
(Proposition~\ref{prop:adequacy}(ii)). Both are configurable per project and neither is
hard-coded. \\
\addlinespace

A10 & Local urgency follows a stationary AR(1)
(Section~\ref{sec:methods_forecast})
& Testable, untested
& The projected urgency path is wrong in shape, and with it every $p_h$. Nineteen turns is
short for a specification test, which is why the coefficient is clipped to $[0,0.98]$ rather
than trusted. \\
\addlinespace

A11 & Weeks are independent Bernoulli trials for the adverse-event rate
(Section~\ref{sec:methods_forecast})
& Testable, untested
& Clustered events, one incident producing three weeks of consequences, break independence
and the Beta posterior understates uncertainty. The prior strength of $N_0 = 26$ dominates
nineteen turns of data, which masks the effect here and would not on a longer campaign. \\
\addlinespace

A12 & The calibration service's outcome definition is the right target
(Section~\ref{sec:methods_scoring})
& Tested
& Every score in this thesis moves. This is the one assumption whose failure was measured
rather than argued: Section~\ref{sec:res_target_trap} scores identical forecasts against
three definitions and finds a spread wider than any modelling decision taken in this work. \\

\end{longtable}

\subsection{What the register says when read as a whole}

Twelve assumptions, of which three are tested, five are testable and untested, one is a
declared modelling choice, and three sit outside what this design can reach. That last group is the honest statement of the thesis's
boundary, and it has a shape: A3 and A4 both concern the declarant, and A6 concerns
correlation between nodes. What the instrument cannot see is people and common causes,
which is to say the two things a supply chain is made of besides its arithmetic.

\begin{figure}[htbp]
\centering
\begin{tikzpicture}[x=1cm, y=1cm, font=\scriptsize,
    st/.style = {minimum width=0.85cm, minimum height=0.85cm, draw=altenNavy,
                 line width=0.6pt, font=\scriptsize\bfseries}]

\node[titreBande, text width=3.2cm] at (2.0, 2.3) {tested};
\node[titreBande, text width=3.6cm] at (6.4, 2.3) {testable, untested};
\node[titreBande, text width=2.4cm] at (10.3, 2.3) {choice};
\node[titreBande, text width=3.4cm] at (13.0, 2.3) {outside this design};

\foreach \n/\x in {A1/1.0, A2/2.0, A12/3.0}
  \node[st, fill=altenAzure!30] at (\x, 1.5) {\n};
\foreach \n/\x in {A5/5.0, A7/6.0, A9/7.0, A10/8.0, A11/9.0}
  \node[st, fill=altenBlueBg] at (\x, 1.5) {\n};
\node[st, fill=white] at (10.3, 1.5) {A8};
\foreach \n/\x in {A3/12.2, A4/13.2, A6/14.2}
  \node[st, fill=altenOchrePale] at (\x, 1.5) {\n};

\node[etiquette, anchor=north, text width=3.0cm] at (2.0, 0.9) {3};
\node[etiquette, anchor=north, text width=3.0cm] at (7.0, 0.9) {5};
\node[etiquette, anchor=north, text width=3.0cm] at (10.3, 0.9) {1};
\node[etiquette, anchor=north, text width=3.0cm] at (13.2, 0.9) {3};

\draw[decorate, decoration={brace, amplitude=4pt, mirror}, draw=altenOchreDark]
     (11.7, 0.55) -- (14.7, 0.55);
\node[etiquette, anchor=north, text width=5.4cm, text=altenOchreDark] at (13.2, 0.25)
     {A3 and A4 concern the declarant, A6 concerns correlation between nodes: people and
      common causes};

\end{tikzpicture}
\caption[The shape of the register]{Twelve assumptions sorted by status. The three on the
right are what bounds the thesis, and they have a shape: the instrument cannot see people, and
it cannot see a cause acting on several nodes at once.}
\label{fig:register_shape}
\end{figure}


That is not a reason to distrust the measurements reported here. It is the reason they are
reported as measurements of a mechanism rather than as estimates of a rate, and the reason
Section~\ref{sec:future_work} asks first for a human campaign rather than for a larger
simulated one. A larger simulated campaign would tighten every interval in
Section~\ref{sec:res_power} and would move none of the three above.
